El sistema de moderación de SocratiCode opera en dos capas. La primera es de naturaleza pedagógica y se implementa en el \textit{system prompt}: la regla 4 (LÍMITE DE DOMINIO) hace que el propio tutor rechace cualquier pregunta ajena a la programación sin necesidad de ningún componente adicional. La segunda es de seguridad activa y se apoya en un segundo LLM que actúa como juez independiente.

En lugar de listas de palabras prohibidas o expresiones regulares, SocratiCode delega el análisis en un modelo de lenguaje dedicado que responde únicamente con \texttt{OK} o \texttt{NO}. Los criterios de bloqueo son estrictos y se limitan a contenido genuinamente dañino (material sexual explícito, instrucciones para el daño físico, discurso de odio o fabricación de armas). El modelo está expresamente instruido para marcar como \texttt{OK} cualquier fragmento relacionado con programación, evitando falsos positivos que bloquearían la actividad educativa. Los modos de moderación (entrada, salida, ambos o ninguno) son configurables en tiempo de ejecución desde el panel de administración.

\paragraph{Moderación de entrada.}

Antes de que el mensaje del alumno llegue al tutor principal, se evalúa de forma síncrona junto con el código del editor. Si el veredicto es negativo, el flujo se interrumpe antes de que el tutor procese nada, el mensaje se marca como \texttt{moderated=True} en base de datos y se devuelve al cliente la respuesta predeterminada:

\PythonCode[lst:input_mod]{Bloqueo de entrada moderada}{Corte del flujo ante un veredicto negativo y emisión de la respuesta predeterminada (\texttt{views.py}).}{python/input_mod.py}{1}{13}{1}

Esta capa opera con política \textit{fail-closed}: si el modelo de moderación lanza una excepción por cualquier motivo, se asume que el contenido es inseguro y se bloquea igualmente.

\paragraph{Moderación de salida.}

Esperar a que el tutor termine de generar para moderar a posteriori anularía la experiencia de \textit{streaming}. En su lugar, la moderación de salida funciona con una ventana deslizante configurable (40 palabras por defecto). Cada vez que el texto acumulado supera ese umbral, se lanza una tarea asíncrona con \texttt{asyncio.create\_task} que evalúa todo el texto generado hasta ese momento sin bloquear la generación ni el envío de \textit{tokens}. Si alguna tarea devuelve \texttt{NO}, la generación se interrumpe en el siguiente ciclo y el mensaje predeterminado reemplaza el contenido parcial que el alumno pudo haber visto durante el \textit{stream}, sin persistir el texto truncado en base de datos. Esta capa opera con política \textit{fail-open}: si una tarea lanza una excepción, se asume que el contenido es seguro para no cortar un \textit{stream} legítimo por un problema de infraestructura.
